\scanpage{131}
% archive foliation: 130
Modules des courbes elliptiques : théorie transcendante

La donnée d'une courbe elliptique \inserted{$E$} sur $\mathbb C$ revient à celle \deleted{à} d'un \inserted{$\mathbb C$} espace vectoriel de rang $1$, savoir $t_E$
\[
\tag{1}
V=t_E\qquad\text{(espace tangent à l'origine)}
\]
et d'un ``réseau''
\[
\tag{2}
H\subset V
\]
(un \deleted{groupe abélien} \inserted{$\mathbb Z$-module} \deleted{discret} \inserted{discret, libre} de rang $2$). La donnée du \deleted{groupe} $\mathbb Z$-module $\unsure{\Pi|H}$ permet de récupérer $V$ en tant qu'espace vectoriel réel \inserted{(de rang $2$)} par
\[
\tag{3}
V\simeq H\otimes_{\mathbb Z}\mathbb R\overset{\text{déf}}=H_{\mathbb R}
\]
et $E$ en tant que groupe de Lie réel, par
\[
\tag{4}
E\simeq V/H\simeq H_{\mathbb R}/H,
\]
tandis que la structure complexe de $E$ est déterminée par la donnée d'une structure d'espace vectoriel complexe sur $V$ par passage aux quotients.

Étant donné un espace vectoriel réel $V$, \inserted{\deleted{de rang $2$}} la donnée d'une structure complexe sur $V$ (i.e. d'un endom $u$ de $V$ de carré $-\mathrm{id}_V$) équivaut aussi à celle d'un \inserted{$\mathbb C$-}sous-espace vectoriel $L$ de $V_{\mathbb C}$
\[
\tag{5}
V_{\mathbb C}=V\otimes_{\mathbb R}\mathbb C
\]
tel que la projection canonique
\[
\tag{6}
\begin{aligned}
V_{\mathbb C}&\longrightarrow V=V_{\mathbb R}\\
z&\longmapsto\tfrac12(z+\bar z)=\Re z
\end{aligned}
\]
\deleted{soit} \inserted{induit} un isomorphisme d'espaces vectoriels réels.
\[
\tag{7}
\deleted{\text{Pour un}}\quad\varphi:L\xrightarrow{\sim}V\quad(\mathbb R\text{-iso})
\]
